\section{Metodologia}
\label{chap:metodologia}

Para o Terceiro Trabalho Computacional, implementamos uma cadeia de processamento composta por três técnicas fundamentais de reconhecimento de padrões, com o intuito de analisar como a redução de dimensionalidade e a flexibilidade das fronteiras de decisão influenciam a acurácia global e a estabilidade numérica. O protocolo experimental utiliza o dataset \textit{Wall-Following Robot Navigation} (24 sensores) submetido a 100 rodadas independentes de treinamento e teste.

\subsection{Redução de Dimensionalidade: PCA via SVD}
A Análise de Componentes Principais (PCA) foi implementada através da decomposição em valores singulares (SVD) da matriz de covariância centralizada. O objetivo é projetar a matriz de dados $\mathbf{X} \in \mathbb{R}^{p \times N}$ em um subespaço $\mathbb{R}^{q}$ com $q < p$, mantendo a máxima variância explicada. 

Diferente de abordagens baseadas puramente na variância acumulada (como a regra dos 95\%), adotamos um critério de otimização empírica para determinar o número ótimo de componentes $q$. Para cada classificador (CQG e DMP), realizamos uma varredura (grid search) de $q \in \{2, 5, 8, 12, 15, 18, 21, 24\}$, executando 100 rodadas de Monte Carlo para cada valor. O valor final de $q$ foi selecionado como aquele que maximiza a acurácia média de validação, permitindo identificar o ponto ideal onde a remoção de ruído compensa a perda de sinal para cada arquitetura específica.

\subsection{Classificador Quadrático Gaussiano (CQG)}
O CQG modela cada classe $w_i$ como uma densidade multivariada $\mathcal{N}(\boldsymbol{\mu}_i, \mathbf{C}_i)$. A função discriminante quadrática utiliza a distância de Mahalanobis e o log-determinante da matriz de covariância específica de cada classe. Para contornar problemas de singularidade, mantivemos a regularização de Tikhonov ($\mathbf{C}_i + 10^{-6} \mathbf{I}$). O CQG permite capturar relações complexas entre sensores, gerando fronteiras de decisão hiper-quadráticas que se adaptam à dispersão dos dados de navegação do robô.

\subsection{Distância Mínima ao Protótipo (DMP) Multi-Protótipo}
O classificador DMP Multi-Protótipo representa uma inovação em relação ao DMP simples. Em vez de um único protótipo por classe, utilizamos o algoritmo K-médias para identificar centros de massa locais. O desafio de determinar o número ótimo de protótipos ($K$) por classe foi resolvido via validação cruzada de clusters, utilizando uma "votação por moda" de seis índices: Davies-Bouldin, Calinski-Harabasz, Dunn, Silhouette, I-Index e Ball-Hall. Esta abordagem permite que o DMP capture a multimodalidade das classes, tornando-o um classificador não-linear potente mesmo utilizando a distância Euclidiana simples para a decisão final.

O protocolo de avaliação comparou o desempenho dos modelos em dois cenários: utilizando os 24 sensores originais e utilizando as componentes principais otimizadas para cada modelo. Esta comparação visa quantificar o impacto da compressão de atributos na capacidade de generalização e no tempo de execução.
